\chapter{Functional Core Languages and Evaluation Semantics}
\label{ch:functional-core-semantics}

A function value contains more than the text of its body.
It must remember which bindings its free variables denote, and its caller must agree on when arguments are evaluated.
Getting either choice wrong changes program meaning before optimization begins.
The functional implementation starts from a small core whose lexical scope and evaluation order are executable contracts.

\section{Choose a core that exposes the hard decisions}
\label{sec:ch89-syntax}
The lab language contains integers, booleans, variables, lambda abstraction, application, let, restricted recursive let, conditionals, integer primitives, and lists with a two-arm case.
It uses S-expressions to keep parsing separate from functional compilation.
For example, \code{(lam x (add x 1))} is a unary function, and \code{(app f 4)} calls a computed function value.
Application is explicit rather than inferred from every parenthesized list.

Table~\ref{tab:ch89-syntax} defines the essential forms.
The implementation validates arities and binder shapes after parsing.
Recursive bindings require a lambda right side, avoiding uninitialized recursive values in the strict machine.
Lists are homogeneous under the type system, but their representation differs between strict and lazy machines.

\begin{table}[htbp]
\caption{The functional core makes binding and evaluation constructs explicit.}
\label{tab:ch89-syntax}
\centering
\begin{tabular}{ll}
\toprule
Form & Meaning \\
\midrule
\code{(lam x body)} & capture lexical environment; bind x on call \\
\code{(app f a)} & apply a computed function \\
\code{(let x value body)} & nonrecursive lexical binding \\
\code{(letrec f lambda body)} & recursive function binding \\
\code{(if test yes no)} & evaluate only the selected arm \\
\code{(cons head tail)} & construct a list cell \\
\code{(case xs empty (h t body))} & branch on list shape; bind fields \\
\bottomrule
\end{tabular}
\end{table}

The syntax is deliberately smaller than ML, Haskell, or Scheme.
There are no modules, mutation, user-defined effects, general user-declared algebraic types, type classes, or exception handlers.
Those features require additional static and runtime machinery rather than mere parser extensions.
The chapters discuss their design implications while the case study states exactly what runs.

\section{Represent lexical scope with environments}
\label{sec:ch89-environment}
An environment $\rho$ maps resolved variable identities to runtime bindings.
A closure is a triple consisting of parameter, body or code pointer, and captured environment.
The evaluation judgment $\rho\vdash e\Downarrow v$ states that expression e evaluates to value v in that environment.
The strict machine evaluates function and argument before entering the closure body with its captured environment extended by the parameter binding.

Consider binding x to 11, binding f to a function adding x to its argument, then shadowing x with 99 before calling f on 4.
Lexical scoping yields 15.
Looking up x in the caller's environment would instead produce 103, an accidental dynamic-scope implementation.
Figure~\ref{fig:ch89-closure} shows the retained binding.

\begin{figure}[htbp]
\centering
\begin{tikzpicture}[node distance=14mm and 16mm]
\node[basic block] (f) {closure f: add captured x to y};
\node[basic block,below left=of f] (old) {captured x = 11};
\node[basic block,below right=of f] (new) {caller shadows x = 99};
\draw[cfg edge] (f)--node[left,font=\small]{uses}(old);
\draw[dashed] (f)--node[right,font=\small]{not captured}(new);
\end{tikzpicture}
\caption{A closure refers to its definition-site binding, not a same-spelled caller binding.}
\label{fig:ch89-closure}
\end{figure}

The Python reference uses validated names and copied environment dictionaries of binding cells.
This correctly models lexical shadowing for the supported syntax but is not a compact production closure layout.
A compiler normally resolves names to declaration identities or lexical addresses once, then stores only needed free variables rather than copying unrelated bindings at every closure construction.
Those representation changes are developed in Chapter~\ref{ch:closures-lifting-combinators}.

\section{State argument evaluation order}
\label{sec:ch89-strictness}
\term{Call-by-value} evaluates the argument before function entry.
Call-by-name delays argument evaluation and can repeat it at each use.
\term{Call-by-need} also delays it but updates shared suspended computation after its first evaluation.
These strategies can agree on terminating pure examples and disagree on termination or resource use.

Let loop be a recursive function that immediately calls itself.
Applying the constant function returning seven to loop(0) diverges under call-by-value and returns seven under call-by-need, because the parameter is unused.
Table~\ref{tab:ch89-strategies} separates the relevant behaviors.
The lab tests this difference explicitly rather than claiming strict and lazy machines implement identical termination semantics.

\begin{table}[htbp]
\caption{Argument strategy changes when work is performed and whether it is repeated.}
\label{tab:ch89-strategies}
\centering
\begin{tabular}{lll}
\toprule
Strategy & Unused divergent argument & Repeated parameter use \\
\midrule
call-by-value & evaluated before entry & reuse argument value \\
call-by-name & not demanded & may recompute expression \\
call-by-need & not demanded & memoize one demanded computation \\
\bottomrule
\end{tabular}
\end{table}

The strict interpreter and stack VM implement left-to-right call-by-value.
The lazy interpreter is an alternate call-by-need semantics over the same surface forms.
Integer primitives demand both operands, a conditional demands only its test and chosen branch, and a list case demands only the outer list shape before selecting an arm.
All machines use mathematical integers through Python's arbitrary-precision arithmetic, not machine-width overflow.

\section{Evaluate a higher-order program step by step}
\label{sec:ch89-evaluate}
Algorithm~\ref{alg:ch89-app} gives strict application.
The closure environment is copied or extended without changing the caller's lexical bindings.
The environment's binding cells allow recursive closures to see a slot installed before the closure is created and filled immediately afterward.
Because recursive right sides are lambdas, closure construction does not read an uninitialized recursive value.

\begin{algorithm}[htbp]
\caption{Evaluate strict application under lexical scope.}
\label{alg:ch89-app}
\begin{pseudocode}
\Require typed application app(f,a), lexical environment $\rho$
\Ensure result under left-to-right call-by-value, or evaluation failure/budget status
\State evaluate $f$ in $\rho$ to closure $(x,body,\rho_c)$
\State evaluate $a$ in $\rho$ to argument value $v$
\State extend $\rho_c$ with a new binding $x\mapsto v$
\State evaluate body in that extended environment
\end{pseudocode}
\end{algorithm}

For the captured-x example, closure creation retains x=11; the later let allocates a different cell for x=99.
Calling f evaluates argument four, binds y=4 in the closure environment, and adds eleven and four.
Dictionary copying costs the environment size in the reference interpreter, so an application is not universally constant-time even before evaluating its body.
Lexical-address arrays or linked frames make different space/time tradeoffs and require their own lookup bounds.

Figure~\ref{fig:ch89-recursion} shows the recursive knot for factorial.
The slot exists before closure construction, and the closure's captured dictionary contains that slot.
Filling the slot afterward establishes the back-edge.
This is a heap graph cycle with a valid closure value, not an attempt to unfold a finite syntax tree cyclically.

\begin{figure}[htbp]
\centering
\begin{tikzpicture}[node distance=25mm]
\node[basic block] (slot) {binding cell fact};
\node[basic block,right=of slot] (closure) {factorial closure};
\draw[cfg edge] (slot.north east) to[bend left=20] (closure.north west);
\draw[cfg edge] (closure.south west) to[bend left=20] node[below,font=\small]{captures cell}(slot.south east);
\end{tikzpicture}
\caption{A recursive function environment ties a cell to the closure that captures it.}
\label{fig:ch89-recursion}
\end{figure}

\section{Separate semantic outcomes from host resource limits}
\label{sec:ch89-budget}
The reference machines charge a finite execution budget and can also encounter Python's recursion or memory limits.
Exhausting fuel is an inconclusive execution outcome, not a mathematical proof that the source diverges.
Algorithm~\ref{alg:ch89-run} keeps parsing, typing, optimization, and execution failures separate from a successful scalar result.
The CLI prints only integer or Boolean results; internal lists and closures are not silently serialized as if they were portable source values.

\begin{algorithm}[htbp]
\caption{Run a functional program through explicit validation and execution boundaries.}
\label{alg:ch89-run}
\begin{pseudocode}
\Require source text, chosen strict or lazy semantics, finite execution budget
\Ensure typed scalar result or explicit failure without a false successful value
\State parse one complete expression and validate shape/binder restrictions
\State infer its type; reject unresolved typing errors
\State apply only optimizations sound for the chosen contract
\State execute with the selected machine and budget
\State publish a scalar only after successful completion
\end{pseudocode}
\end{algorithm}

Parsing the finite token stream costs linear scanning plus tree construction, subject to token and nesting limits.
Type inference and execution have separate costs; general recursive source evaluation is not bounded by source size.
Optimization cannot erase a potentially divergent strict argument merely because it returns no useful numeric contribution.
These distinctions are central to testing transformations on functional code.

\section{Relate environment execution to substitution semantics}
\label{sec:ch89-substitution}
An alternative account of pure lambda evaluation substitutes an argument into the function body.
That substitution must avoid capturing free variables under nested binders.
The environment machine represents the same binding relationships indirectly: the body retains its lexical references and the environment supplies values for them.
Correctness of closure execution can be argued by relating an environment to a capture-avoiding substitution for its bound declarations.

At function creation, the captured environment records the substitutions justified at the definition site.
At application, extending that captured environment by the parameter value corresponds to substituting the argument for that particular binder.
Shadowing creates a fresh declaration relationship rather than overwriting the meaning of an outer captured reference.
The captured-x example tests exactly that distinction.
An implementation that executes the body in the caller environment breaks the relation and exhibits the dynamic-scope result 103 instead of fifteen.

\section{Exercises}
\label{sec:ch89-exercises}
\begin{enumerate}
\item\label{ex:ch89-scope} Evaluate the captured-x example and contrast the erroneous dynamic-scope result.
\item\label{ex:ch89-strict} Compare strict and lazy evaluation of a constant function applied to loop(0).
\item\label{ex:ch89-recursion} Explain why allocating the recursive binding cell before constructing a closure is safe for a lambda right side but not arbitrary eager self-reference.
\item\label{ex:ch89-budget} A recursive program consumes all fuel.
What does this establish, and why must the CLI not print an arbitrary partial scalar?
\end{enumerate}

\section{Further reading}
\label{sec:ch89-reading}
Appel covers lexical environments and closures~\cite{appel-1998-modern}.
The STG paper develops compiled non-strict execution~\cite{peytonjones-1992-stg}.
